% =====================================================================
% Appendix I --- From the input--output table to Omega: the data pipeline
% ---------------------------------------------------------------------
% Job: document, step by step, how the published Destatis input--output
%      table becomes the calibrated objects the kernel consumes --- above
%      all the input-share matrix Omega, plus alpha, lambda and the
%      consumption shares omega. Section 4 states the definitions
%      (eq:io-row, eq:io-column); Appendix D carries the A-bill leak
%      accounting; Appendix F the investment vector. This appendix fills the
%      gap between "the table" and "the shares".
% Code: OpenCGE@src/core/accounting.jl (read_data, final_demand_split,
%      generate_data, conditional_input_shares, assemble_data) and
%      OpenCGE@src/core/calibration.jl (retained_dataset for variants).
% Status: STUB --- placeholder outline (2026-10-12).
% =====================================================================
\section{From the input--output table to $\Omega$: the data pipeline}\label{app:pipeline}

% TODO: short opening paragraph: what goes in (one published table), what
% comes out (Omega, alpha / factor share, lambda, omega, the leak vectors),
% and why it matters --- every closure in the paper is calibrated on the
% same objects, so differences in results cannot come from the data step.
% Possibly a small flow diagram (TikZ): raw table -> cleaned table ->
% Z, Z_dom -> Omega / alpha / lambda / omega -> kernel.

\subsection{Source data}\label{app:pipeline-source}
% TODO: the Destatis input--output table for Germany 2019 (revision 2019),
% 71 product groups, product-by-product, basic prices; file name in data/
% (destatis_input_output_2019_revision2019.xlsx -> I-O_DE2019_formatiert.csv);
% citation of the Destatis release (cross-reference app:provenance).
% TODO: layout of the table as read: rows 1--71 the sector block, then the
% valuation rows (row 73 domestic use, row 74 imports, row 75 product taxes
% less subsidies, row 76 their sum), the value-added rows and gross output;
% columns 1--71 the users, then the final-demand categories.
% TODO: the formatting step from xlsx to csv (what is removed or renamed,
% decimal/missing-value conventions "-" and "x") --- state it is mechanical.

\subsection{The intermediate-use matrix}\label{app:pipeline-z}
% TODO: Z = intermediate deliveries x_{ui} (supplier u to user i), domestic
% and imported combined; link to eq:io-row / eq:io-column.
% TODO: imports are reported by USING sector only (row 74), not by supplier:
% the proportional allocation assumption that yields the domestic matrix
% Z_dom (each user keeps its supplier composition).
% TODO: the clamping guard --- the few sectors whose recorded intermediate
% imports exceed intermediate use; list them, give their size, and say they
% are treated as fully import-dependent in Z_dom.

\subsection{Normalizing to cost shares: $\Omega$}\label{app:pipeline-omega}
% TODO: Omega_{iu} = p_u x_{ui} / sum_{u'} p_{u'} x_{u'i}: column of Z
% divided by the user's total intermediate bill, transposed to user x
% supplier; each row sums to one (the contract the CES price index
% eq:model-ip relies on: unit cost = 1 at the calibration point).
% TODO: which Omega the kernel uses --- the raw (domestic + imported) Omega
% for equilibrium technology vs. the domestic Omega_dom kept for incidence
% and audit; justify the choice and point to the A-bill accounting in
% app:kernel (the domestic bill A_bill charged to intermediates, M_int and
% T_int booked as leaks).
% TODO: descriptive statistics on Omega: sparsity, diagonal (self-supply)
% shares, the high-self-share sectors (> 0.45) that motivate the "reduced"
% variant (app:variants); maybe a heat map of Omega as a figure.

\subsection{Value added, factor shares and Domar weights}\label{app:pipeline-va}
% TODO: value added decomposed into compensation of employees, other
% production taxes less subsidies, depreciation, net operating surplus
% (identity checked to 1e-9); the composite factor share alpha_i = v_i /
% (p_i y_i) --- state clearly that this is total value added, not pure
% labour, and why the kernel uses it.
% TODO: Domar weights lambda_i = gross output / GDP and the labour-share
% weights lambda_i * alpha_i (summing to one).

\subsection{Final demand and household shares $\omega$}\label{app:pipeline-fd}
% TODO: the final-demand categories (household consumption, government,
% investment, exports); purchaser-price categories split proportionally into
% domestic, imported and product-tax components (final_demand_split; source
% limitation: imports and taxes on final use are reported by category).
% TODO: consumption shares omega calibrated so the model clears at y =
% lambda (from (I - diag(1 - alpha) Omega)' y); negative entries clamped to
% zero and renormalized --- report how many and how large.
% TODO: government, investment and export blocks and the implied saving rate
% s (data-implied residual, not assumed); cross-reference app:kernel.

\subsection{Consistency checks}\label{app:pipeline-checks}
% TODO: the three GDP sides (production = income exactly; production vs.
% expenditure residual 5.387 %, an open item of the raw table itself) and
% the 10 % guard.
% TODO: sum_u Omega_{iu} = 1 for every i; sum lambda_i alpha_i = 1;
% sum A + sum M_int + sum T_int = sum lambda - 1 to 2e-16 (app:kernel).
% TODO: the paper-side tests that assert these on the real table
% (tests/test_calibration.jl) and the 1e-8 reproduction baseline.

\subsection{Sector-exclusion variants}\label{app:pipeline-variants}
% TODO: how the "70s" and "reduced" variants are rebuilt from the RAW table
% (retained_dataset) rather than by slicing the full-table matrices: Omega
% rows renormalized on retained suppliers, lambda and alpha re-derived,
% final-demand imports/taxes scaled by the retained share. Keep short and
% point to app:variants for the results.
